\documentclass[11pt,a4paper]{ctexart}
\usepackage[margin=25mm]{geometry}
\usepackage{amsmath,amssymb,amsthm,mathtools,booktabs,tabularx}
\usepackage{xcolor,hyperref,listings,fancyhdr,lastpage}
\setlength{\parindent}{2em}
\setlength{\parskip}{0.24em}
\definecolor{dblue}{HTML}{17364B}
\definecolor{dteal}{HTML}{176F78}
\definecolor{codebg}{HTML}{F4F6F8}
\hypersetup{colorlinks=true,linkcolor=dteal,urlcolor=dteal}
\pagestyle{fancy}\fancyhf{}
\fancyhead[L]{\small 异质 $K=5$ 稳健投票：候选闭式与完整性证明路线}
\fancyhead[R]{\small D2 阶段更新}
\fancyfoot[C]{\small \thepage\ / \pageref{LastPage}}
\newtheorem{theorem}{定理}[section]
\newtheorem{lemma}[theorem]{引理}
\newtheorem{proposition}[theorem]{命题}
\newtheorem{corollary}[theorem]{推论}
\theoremstyle{definition}
\newtheorem{remark}[theorem]{说明}
\newenvironment{proofcn}{\par\noindent\textbf{证明.}\ }{\hfill$\square$\par}
\lstdefinestyle{pycode}{language=Python,basicstyle=\ttfamily\scriptsize,backgroundcolor=\color{codebg},frame=single,rulecolor=\color{gray!35},columns=fullflexible,breaklines=true,breakatwhitespace=false,showstringspaces=false,keepspaces=true,tabsize=4,numbers=left,numberstyle=\tiny\color{gray},numbersep=5pt}
\lstset{style=pycode}
\newcommand{\Risk}{\mathsf R}
\newcommand{\calP}{\mathcal P}
\newcommand{\calC}{\mathcal C}
\newcommand{\ind}{\mathbf 1}
\newcommand{\EE}{\mathbb E}
\newcommand{\PP}{\mathbb P}
\newcommand{\comp}[1]{\overline{#1}}
\title{\bfseries 异质五分类器稳健投票：\\[-2pt]98 项证书候选的不可删减性与有限完整性判据}
\author{D2 理论研究工作稿}
\date{2026 年 10 月 13 日}
\begin{document}
\maketitle
\begin{abstract}
研究已知边际错误率 $q_i$ 和共同出错概率上界 $b_{ij}$ 的五个二元分类器，目标是在全部确定性标签对称（self-dual）投票规则上最小化最坏错误率。此前的简单候选 $\min\{q_{\min},c^*,B_5\}$ 已被反例否定。本文只保留当前最主要的成果：给出\textbf{98 项、8 类公式结构}组成的显式仿射上界 $U_{98}$；严格证明全部 98 项均为合法的二次矩证书；构造\textbf{9 组精确有理数见证参数}，通过模型置换证明每一项都在区域 $R$ 内不可删除；并给出以参数分区顶点和\textbf{60 类阈值规则}为基础的候选完整性\emph{充要判据}。本文\textbf{尚未证明} $\Risk_{\rm det}^*=U_{98}$，也未宣称 98 项是全体合法证书中的最少项数。附录包含所有用于构造证书、核查不可删减性和核查规则同构数目的完整 Python 代码。
\end{abstract}

\section{问题与研究状态}
设 $X_i\in\{0,1\}$ 表示第 $i$ 个二元分类器是否出错，$q_i=\PP(X_i=1)\in(0,1/2)$，已知上界 $\EE[X_iX_j]\le b_{ij}$，其中 $0\le b_{ij}\le\min\{q_i,q_j\}$。令 $x\in\{0,1\}^5$，定义模糊集
\begin{equation}\label{eq:ambig}
\calP(q,b)=\left\{\pi\ge0: \sum_x\pi_x=1,\ \sum_xx_i\pi_x=q_i,\ \sum_xx_ix_j\pi_x\le b_{ij}\ \forall i<j\right\}.
\end{equation}
始终假定 $\calP(q,b)\ne\emptyset$。标签对称的确定性错误规则 $F$ 满足 $F(x)+F(\mathbf1-x)=1$。其最坏错误率与全局最优值为
\begin{equation}\label{eq:game}
\Risk(F;q,b)=\max_{\pi\in\calP(q,b)}\sum_xF(x)\pi_x,\qquad
\Risk_{\rm det}^*(q,b)=\min_{F\text{ self-dual}}\Risk(F;q,b).
\end{equation}
研究域为
\begin{equation}\label{eq:R}
R=\left\{(q,b):\sum_iq_i-3q_{\min}\le1,\ \calP(q,b)\ne\varnothing\right\},\qquad q_{\min}=\min_iq_i.
\end{equation}
早期三项猜想 $\Risk_{\rm det}^*=\min\{q_{\min},c^*,1-\sum q_i+\sum b_{ij}\}$ 在 $R$ 内\textbf{不成立}。当前 98 项候选为一个已证\emph{上界}，而不是已证闭式：
\begin{equation}\label{eq:ub}
\Risk_{\rm det}^*(q,b)\le U_{98}(q,b),\qquad
\Risk_{\rm det}^*(q,b)\stackrel{?}{=}U_{98}(q,b)\quad((q,b)\in R).
\end{equation}

\section{98 项显式二次证书}
\subsection{总公式}
对任意二次多项式
\begin{equation}\label{eq:quadratic}
Q(x)=a+\sum_{i=1}^{5}\beta_ix_i+\sum_{1\le i<j\le5}\gamma_{ij}x_ix_j,\qquad \gamma_{ij}\ge0,
\end{equation}
若 $Q(x)\ge F(x)$ 对所有 $x$ 成立，则
\begin{equation}\label{eq:upper}
\Risk(F;q,b)\le L_Q(q,b):=a+\sum_i\beta_iq_i+\sum_{i<j}\gamma_{ij}b_{ij}.
\end{equation}
这是逐点取期望以及 $\gamma_{ij}\ge0$ 的直接推论。令 $\calC_{98}$ 是下表指定的 98 个二次证书，定义
\begin{equation}\label{eq:U}
\boxed{U_{98}(q,b):=\min_{Q\in\calC_{98}}L_Q(q,b).}
\end{equation}
表格中的“模板数”表示模型编号重新排列后产生的不同仿射式数量。
\begin{center}
\begin{tabularx}{\textwidth}{@{}p{.46\textwidth}Xp{18mm}@{}}
\toprule
家族 & 代表性规则或结构 & 项数\\\midrule
单分类器（独裁） & $w=(1,0,0,0,0)$ & 5\\
奇环并集证书 & 所有三角形及五环 & 22\\
加权平移二次矩证书 & $w=(3,1,1,1,1)$ & 5\\
加权平移二次矩证书 & $w=(2,1,1,1,0)$ & 20\\
加权平移二次矩证书 & $w=(2,2,1,1,1)$ & 10\\
加权平移二次矩证书 & $w=(1,1,1,1,1)$ & 1\\
根节点修正证书 & 一个特殊模型 $i$ & 5\\
三核心修正证书 & $w=(3,2,2,1,1)$ & 30\\\midrule
\textbf{合计} & \textbf{8 类代数结构；9 个置换轨道（区分三环、五环）} & \textbf{98}\\\bottomrule
\end{tabularx}
\end{center}

\subsection{权重通用的平移二阶矩证书}
取非负整数权重 $w$，总和 $\sum_iw_i=2r-1\ge3$。记 $N_w=\sum_iw_iX_i$。取多项式
\begin{equation}\label{eq:BP}
H_r(t)=\frac{(t-r+2)(t-r+1)}2.
\end{equation}
对任意整数 $t\ge0$，两个因子是连续整数，因此 $H_r(t)\ge0$；当 $t\ge r$ 时 $H_r(t)\ge1$。于是 $H_r(N_w)$ 逐点控制权重 $w$ 的多数失败指示函数。利用 $X_i^2=X_i$ 展开：
\begin{equation}\label{eq:BPbound}
\boxed{\begin{aligned}
B_w={}&\frac{(r-2)(r-1)}2
+\frac12\sum_i w_i(w_i-2r+3)q_i\\
&+\sum_{i<j}w_iw_jb_{ij}.
\end{aligned}}
\end{equation}
此处的具体逐点证书可独立证明；它与 Boros--Pr\'ekopa 的二项矩概率界属于同一方法脉络\cite{BP1989}，不能据此声称该经典界是本研究原创。

\subsection{根节点与三核心修正}
取根节点 $i$，对应 $w_i=3$、其他权重为 $1$。令 $T=\sum_{j\ne i}x_j$。根节点二次证书为
\begin{equation}\label{eq:root}
Q_{\rm root,i}(x)=\frac13\left(1-\sum_{j=1}^5x_j+\sum_{\substack{j<k\\j,k\ne i}}x_jx_k+4\sum_{j\ne i}x_ix_j\right).
\end{equation}
若 $x_i=0$，其值为 $(T-1)(T-2)/6$，在 $T=0,1,2,3,4$ 时非负，且唯一失败状态 $T=4$ 对应值 $1$；若 $x_i=1$，其值为 $T+\binom{T}{2}/3$，对每个失败状态 $T\ge1$ 至少为 $1$。因此~\eqref{eq:root} 为合法证书，给出上界
\begin{equation}\label{eq:rootbound}
B_{\rm root}(i)=\frac{1-\sum_jq_j+\sum_{\substack{j<k\\j,k\ne i}}b_{jk}+4\sum_{j\ne i}b_{ij}}3.
\end{equation}
对 $w$ 为 $(3,2,2,1,1)$ 的排列，记 $H=\{i:w_i\ge2\}$，$N_w=\sum_iw_ix_i$，$T_H=\sum_{i\in H}x_i$。定义
\begin{equation}\label{eq:core}
Q_{\rm core}(x)=\frac{(N_w-3)(N_w-4)}2-\left(1-T_H+\binom{T_H}{2}\right).
\end{equation}
当 $T_H=1$ 或 $2$ 时，括号为 $0$，证书仍是原本合法的 $H_5(N_w)$；当 $T_H=0$ 时 $N_w\in\{0,1,2\}$，故 $H_5(N_w)\ge1$，减去 $1$ 后仍非负，而投票不会失败；当 $T_H=3$ 时 $N_w\ge7$，有 $H_5(N_w)\ge6$，减去 $1$ 后仍大于 $1$。所以~\eqref{eq:core} 同样逐点上界失败指示函数。展开它即可得到 30 条显式仿射证书，并且每条二次系数 $\gamma_{ij}$ 非负。

奇环证书取 $Q_C(x)=\sum_{ij\in C}x_ix_j$，因为奇环上的多数错误必包含一对相邻错误顶点。单模型证书取 $Q_i=x_i$。以上逐点论证证明了~\eqref{eq:ub} 的第一个不等式。完整系数生成算法见附录 A。

\section{新证明：98 个仿射项在 $R$ 内不可删减}
\begin{theorem}[候选族内部的不可删减性；精确有理数认证]\label{thm:irred}
对每个 $Q\in\calC_{98}$，存在满足 $0<q_i<1/2$、$\calP(q,b)\ne\varnothing$、$(q,b)\in R$ 的参数，使
\begin{equation}\label{eq:strict}
L_Q(q,b)<\min_{Q'\in\calC_{98}\setminus\{Q\}}L_{Q'}(q,b).
\end{equation}
因而删除 98 项中的任何一项，都将使 $U_{98}$ 在某个可行的 $R$ 内参数处严格变大。
\end{theorem}
\begin{proofcn}
所有证书在对 $5$ 个模型编号同步置换后封闭。按这一群作用，98 项分成 9 个置换轨道，数量依次为
\[
5,\;10,\;12,\;5,\;20,\;10,\;1,\;5,\;30,
\]
总计 98。只需为每个轨道提供一个见证参数。

附录 B 用精确有理数列出 9 个错误模式联合分布及各对上界的原始见证。为保证 $q_i$ 严格位于 $(0,1/2)$，将原见证以系数 $99/100$ 和五个独立 Bernoulli$(1/5)$ 错误分布以系数 $1/100$ 混合；同时以相同系数混合共同出错上界（参考上界为 $1/25$）。凸组合仍可行，且满足 $R$。脚本对每个见证\emph{逐项}用 \texttt{fractions.Fraction} 检查：联合分布归一化与非负性、五个边际、10 条成对上下界、$R$ 条件，以及所选项相对其他 97 项的\emph{严格}差值。所得到的最小差值见表~\ref{tab:witness}；它们全部为严格正的有理数。因为模型置换不会改变全部约束及严格差值，每个轨道的全部项都存在见证，定理得证。这是可逐行复核的\textbf{有限计算机辅助证明}。
\end{proofcn}
\begin{table}[ht]
\centering\small
\caption{9 个代表证书与其在精确见证参数下胜过其余 97 项的最小严格差值。}\label{tab:witness}
\begin{tabular}{@{}llr@{}}
\toprule
家族（代表标签） & 代表权重/结构 & 最小差值 $\Delta>0$\\\midrule
\texttt{dict:0} & 独裁 & $2467/10000$\\
\texttt{cycle:0} & 三角形 & $33/100$\\
\texttt{cycle:10} & 五环 & $247/625$\\
\texttt{bp:11113} & $(3,1,1,1,1)$ & $1389/10000$\\
\texttt{bp:01112} & $(2,1,1,1,0)$ & $539/2500$\\
\texttt{bp:11122} & $(2,2,1,1,1)$ & $116/625$\\
\texttt{bp:11111} & 五人普通多数 & $4399/20000$\\
\texttt{rooted:0} & 根节点修正 & $1814/24375$\\
\texttt{corrected:11223} & 三核心修正 & $737/5000$\\
\bottomrule
\end{tabular}
\end{table}
\begin{remark}[这一定理\emph{没有}证明什么]
“98 项不可删”只表示\emph{固定的这 98 个仿射式}之间没有冗余；如果引入一个不同的、具有更强上界作用的新证书，仍然可能得到项数更少的另一组公式。更重要的是，该结果绝不意味着所有合法二次证书已经被这 98 项穷尽。
\end{remark}

\section{推进完整性证明：参数分区顶点充要判据}
\subsection{固定规则最坏错误率的凹性}
对固定 self-dual 规则 $F$，其最坏联合分布 LP 的对偶为
\begin{equation}\label{eq:fixeddual}
\Risk(F;q,b)=\min_{\substack{a,\beta\text{ 自由},\;\gamma\ge0\\a+\sum_i\beta_ix_i+\sum_{i<j}\gamma_{ij}x_ix_j\ge F(x)\ \forall x}}
\left(a+\sum_iq_i\beta_i+\sum_{i<j}b_{ij}\gamma_{ij}\right).
\end{equation}
LP 强对偶在 $\calP(q,b)$ 非空时成立。因对偶可行域不依赖参数 $(q,b)$，式~\eqref{eq:fixeddual} 是仿射函数族的下确界，故固定规则值函数 $\Risk(F;q,b)$ 关于 $\theta=(q,b)$ 凹。

\begin{lemma}[最优确定性规则可选择线性阈值]
若 $\calP(q,b)$ 非空，则存在最优确定性 self-dual 规则可由某组实数权重 $w$ 的线性符号投票实现。
\end{lemma}
\begin{proofcn}
取实现全局 minimax 值的规则 $F$，并取其固定规则 LP 的最优二次对偶证书 $Q\ge F$。由 self-dual 性质，每对互补模式 $x,\mathbf1-x$ 至少一侧 $Q\ge1$。将 $Q$ 较大的那侧判为失败，便得到另一条合法 self-dual 规则 $F_Q\le Q$，其最坏错误率不高于原规则。又由于
\[
Q(x)-Q(\mathbf1-x)
=\sum_{i=1}^5\left(\beta_i+\frac12\sum_{j\ne i}\gamma_{ij}\right)(2x_i-1),
\]
决定选哪一侧的比较恰是一条带符号线性阈值规则。对双方 $Q$ 值相等的有限个模式，可用任意小的一般位置线性扰动打破平局；原本严格的方向不变，而平局对的两侧都满足 $Q\ge1$，所以仍有 $F_Q\le Q$。这就给出最优线性阈值规则。
\end{proofcn}

\subsection{证书活跃区域划分}
令 $\mathcal M$ 为满足 $0\le q_i\le1/2$、$0\le b_{ij}\le\min(q_i,q_j)$ 以及\eqref{eq:ambig} 非空的闭参数多面体。这是由 $32$ 个 $\pi$ 变量与线性约束定义的多面体的有界投影，因此确实是有理多面体。

为处理 $q_{\min}$，对 $i\in\{1,\dots,5\}$ 定义线性分支
\[
R_i=\{\theta\in\mathcal M:q_i\le q_j\ \forall j,\;\sum_jq_j-3q_i\le1\}.
\]
则闭区域 $\overline R=\bigcup_iR_i$。对每一项 $C\in\calC_{98}$，定义它成为候选最优项的参数区域
\begin{equation}\label{eq:cell}
P_{i,C}=R_i\cap\{\theta:L_C(\theta)\le L_D(\theta)\ \forall D\in\calC_{98}\}.
\end{equation}
每个非空 $P_{i,C}$ 都是\textbf{有界有理多面体}。在其内部，$U_{98}(\theta)=L_C(\theta)$。

\begin{theorem}[完整性可化为有限个参数顶点上的对手证书]\label{thm:vertex}
以下两条陈述等价：
\begin{enumerate}
\item $\Risk^*_{\rm det}(\theta)=U_{98}(\theta)$ 对所有 $\theta\in\overline R$ 成立；
\item 对每个确定性 self-dual 阈值规则 $F$、每个非空分区 $P_{i,C}$、每个顶点 $v\in\operatorname{vert}(P_{i,C})$，存在 $\pi^{F,i,C,v}\in\calP(v)$，满足
\begin{equation}\label{eq:primalproof}
\sum_x F(x)\pi_x^{F,i,C,v}\ge L_C(v).
\end{equation}
\end{enumerate}
其中只需检查 $1{,}882$ 个不同的带符号阈值规则；按模型置换可进一步归成 $60$ 个同构规则代表，其余分区与见证同步置换。
\end{theorem}
\begin{proofcn}
(1)$\Rightarrow$(2)：在任意分区顶点 $v$，有 $\Risk(F;v)\ge \Risk^*_{\rm det}(v)=U_{98}(v)=L_C(v)$。有限维对手 LP 在 $\calP(v)$ 上取到最大值，故存在满足~\eqref{eq:primalproof} 的 $\pi$。所有参数和 LP 系数均为有理数时，可选有理数基可行解。

(2)$\Rightarrow$(1)：固定任意 $\theta\in P_{i,C}$ 与规则 $F$。因为该分区是有界多面体，可将 $\theta$ 写成顶点的凸组合 $\theta=\sum_v\lambda_vv$，$\lambda_v\ge0$、$\sum_v\lambda_v=1$。由固定规则值函数的凹性与~\eqref{eq:primalproof}，
\[
\Risk(F;\theta)\ge\sum_v\lambda_v\Risk(F;v)
\ge\sum_v\lambda_vL_C(v)=L_C(\theta)=U_{98}(\theta).
\]
对全部 self-dual 阈值规则取最小，得到 $\Risk_{\rm det}^*(\theta)\ge U_{98}(\theta)$；反向不等式由~\eqref{eq:ub} 已知，故取等号。原先证明的“二次对偶互补差为线性函数”表明总存在最优的带符号阈值规则，因此无需枚举全部 $65{,}536$ 个 self-dual 规则。最后，整数穷举证明带符号阈值规则只有 $1{,}882$ 个，在 $S_5$ 重标号下只有 $60$ 个轨道，见附录 C。
\end{proofcn}

\begin{remark}[证明工作的实际边界]
定理~\ref{thm:vertex} 是\textbf{真正的等价判据}，不是 $U_{98}$ 已经完整的证明。要确立该猜想，仍需生成各分区\emph{全部}顶点，并在全部必要阈值规则上提供满足~\eqref{eq:primalproof} 的有理数分布。此法把\emph{对偶极点穷举}替换为\emph{参数分区顶点及原始对手证书穷举}；哪一种实际更省计算尚待评估。有限随机测试或部分对偶证书测试都不能替代此步。
\end{remark}

\section{研究状态与下一步}
目前已完成：$(i)$ 98 个合法且显式的仿射上界，$(ii)$ 9 个有理数见证证明全部 98 项在 $R$ 内不可删减，$(iii)$ 以参数分区顶点与 $60$ 类规则为基础的精确完整性充要判据。\textbf{尚未完成}：所有 $P_{i,C}$ 的顶点生成，以及逐个顶点的完整最坏分布证书审计。因而将 $U_{98}$ 直接写成 $\Risk_{\rm det}^*$ 的等式，仍会超出目前已证明的内容。

下一阶段应先利用 $S_5$ 对称性缩减非空参数分区，比较\emph{原始顶点穷举}与\emph{对偶顶点穷举}的复杂度，再选择其中一条执行精确有理数审计。鉴于 98 项在本族内部全部必要，若以\emph{项数更少}为目标，必须先找到一种具有不同系数结构、能同时替换数个旧证书的新构造；单纯剔除旧项行不通。

\begin{thebibliography}{1}
\bibitem{BP1989} E.~Boros and A.~Pr\'ekopa, ``Closed Form Two-Sided Bounds for Probabilities that At Least $r$ and Exactly $r$ Out of $n$ Events Occur,'' \emph{Mathematics of Operations Research}, 14(2):317--342, 1989. \href{https://doi.org/10.1287/moor.14.2.317}{doi:10.1287/moor.14.2.317}.
\end{thebibliography}

\appendix
\section{98 个证书的完整生成及逐点检查代码}
\noindent 用法：保存为 \texttt{D2\_K5\_98\_term\_closed\_form\_candidate.py} 并直接运行。只用标准库和 NumPy。每个候选二次证书均在全部 32 个二元错误状态上被精确有理数不等式核查；此步骤证明候选\emph{有效性}，不是\emph{完整性}。
\begin{lstlisting}
"""D2, K=5: an explicit 98-affine-term robust voting certificate family.

THIS IS A CANDIDATE FOR AN EXACT MINIMAX CLOSED FORM, NOT A PROVED EQUALITY.
Each of its 98 terms is rigorously a valid (possibly non-sharp) worst-case
error upper bound, under feasible marginal / pairwise joint-error caps.

The terms are: 5 dictators + 22 odd cycles + 36 weighted-BP polynomials
+ 5 rooted fractional polynomials + 30 corrected heavy-triplet polynomials.

No LP is required to evaluate this explicit bound. Validity checks evaluate
every polynomial on all 32 binary error patterns.

Requires only Python and numpy (for checks). Compatible with Python>=3.10.
"""
from __future__ import annotations
import itertools
from dataclasses import dataclass
from fractions import Fraction
import numpy as np

K=5
EDGES=tuple(itertools.combinations(range(K),2))
STATES=np.array([[(mask>>i)&1 for i in range(K)] for mask in range(32)],dtype=int)
PAIR=np.array([STATES[:,i]*STATES[:,j] for i,j in EDGES]).T
MOMENT=np.column_stack([np.ones(32,dtype=int),STATES,PAIR])

@dataclass(frozen=True)
class Certificate:
    label: str
    a: Fraction
    beta: tuple[Fraction,...]
    gamma: tuple[Fraction,...]
    weights: tuple[int,...]

    def coefficients(self): return (self.a,)+self.beta+self.gamma
    def bound(self,q,b):
        return self.a + sum(v*Fraction(str(x)) for v,x in zip(self.beta,q)) + \
               sum(v*Fraction(str(x)) for v,x in zip(self.gamma,b))
    def check(self):
        assert len(self.beta)==5 and len(self.gamma)==10
        assert all(z>=0 for z in self.gamma)
        assert sum(self.weights)%2==1 and sum(self.weights)>0
        for x in STATES:
            L=self.a + sum(c*int(t) for c,t in zip(self.beta,x)) + \
                       sum(c*int(x[i])*int(x[j]) for c,(i,j) in zip(self.gamma,EDGES))
            F=int(2*sum(int(v)*int(t) for v,t in zip(self.weights,x))>sum(self.weights))
            assert L>=F,(self.label,tuple(x),L,F)


def odd_cycles():
    cc=[]
    for vertices in itertools.combinations(range(K),3):
        cc.append(tuple(EDGES.index(e) for e in itertools.combinations(vertices,2)))
    for rest in itertools.permutations(range(1,K)):
        cyc=(0,)+rest
        if cyc[1]>cyc[-1]:continue
        cc.append(tuple(EDGES.index(tuple(sorted((cyc[i],cyc[(i+1)%K])))) for i in range(K)))
    assert len(cc)==22
    return cc


def weighted_bp(weights,label):
    w=tuple(weights);total=sum(w)
    assert total>=3 and total%2==1
    r=(total+1)//2
    a=Fraction((r-2)*(r-1),2)
    beta=tuple(Fraction(ww*(ww-2*r+3),2) for ww in w)
    gamma=tuple(Fraction(w[i]*w[j]) for i,j in EDGES)
    return Certificate(label,a,beta,gamma,w)


def rooted(i):
    w=tuple(3 if j==i else 1 for j in range(K))
    a=Fraction(1,3)
    beta=(Fraction(-1,3),)*K
    gamma=tuple(Fraction(4 if i in (j,k) else 1,3) for j,k in EDGES)
    return Certificate('rooted:%d'%i,a,beta,gamma,w)


def corrected_core(weights):
    w=tuple(weights)
    assert sorted(w)==[1,1,2,2,3]
    root=weighted_bp(w,'unused')
    H={i for i,wi in enumerate(w) if wi>=2}
    beta=tuple(root.beta[i]+(1 if i in H else 0) for i in range(K))
    gamma=tuple(root.gamma[e]-(1 if i in H and j in H else 0)
                for e,(i,j) in enumerate(EDGES))
    return Certificate('corrected:%s'%(''.join(map(str,w))),root.a-1,beta,gamma,w)


def make_certificates():
    out=[]
    for i in range(K):
        w=tuple(int(j==i) for j in range(K))
        out.append(Certificate('dict:%d'%i,Fraction(0),tuple(Fraction(v) for v in w),(Fraction(0),)*10,w))
    for ix,C in enumerate(odd_cycles()):
        # Odd cycle votes are a committee of 3 or 5, irrespective of which cyclic edges bound its error.
        vertices={i for e in C for i in EDGES[e]}
        w=tuple(int(i in vertices) for i in range(K))
        gamma=tuple(Fraction(int(e in C)) for e in range(10))
        out.append(Certificate('cycle:%d'%ix,Fraction(0),(Fraction(0),)*5,gamma,w))
    for canonical in ((3,1,1,1,1),(2,1,1,1,0),(2,2,1,1,1),(1,1,1,1,1)):
        for w in sorted(set(itertools.permutations(canonical))):
            out.append(weighted_bp(w,'bp:%s'%(''.join(map(str,w)))))
    for i in range(K):out.append(rooted(i))
    for w in sorted(set(itertools.permutations((3,2,2,1,1)))):
        out.append(corrected_core(w))
    assert len(out)==98,len(out)
    assert len(set(c.coefficients() for c in out))==98
    return out


def bound(q,b):
    candidates=[(c.bound(q,b),c.label) for c in make_certificates()]
    return min(candidates)

if __name__=='__main__':
    cert=make_certificates()
    for c in cert:c.check()
    print('Verified all 98 polynomial majorants exactly on all 32 patterns.')
    q=[Fraction(9,26)]*4+[Fraction(5,26)]
    b=[Fraction(3,26) if 4 not in e else Fraction(0) for e in EDGES]
    v,label=bound(q,b)
    assert v==Fraction(1,26)
    print('Targeted example: bound=',v,'label=',label)
\end{lstlisting}
\section{9 组有理数见证：98 项不可删减性验证代码}
\noindent 保存为 \texttt{D2\_K5\_98\_irredundancy\_exact.py}，与附录 A 的脚本置于同一目录。脚本只用有理数运算逐个检查 9 个代表证书的严格优势，通过模型置换覆盖全部 98 项。

\begin{lstlisting}
"""Exact rational, self-contained irredundancy witnesses for D2's 98 affine certificates.
Run in same directory as D2_K5_98_term_closed_form_candidate.py.
Proves none of these 98 terms can be deleted WITHOUT changing U_98 on
some feasible high-precision R parameter point. It does NOT prove U_98=R_det.
"""
from fractions import Fraction as F
from D2_K5_98_term_closed_form_candidate import EDGES, STATES, make_certificates

# For each of 9 S_5-orbits: (label, rational mass by binary mask, rational b by
# the fixed order EDGES=(01,02,03,04,12,13,14,23,24,34)).
WITNESSES = [
 ('dict:0', {2:'1/4',4:'1/4',8:'1/4',16:'1/4'},
  ['0']*4 + ['1/4']*6),
 ('cycle:0', {2:'1/6',4:'1/6',9:'1/6',10:'1/6',17:'1/6',20:'1/6'},
  ['0','0','1/6','1/6','0','1/6','1/6','1/6','1/6','0']),
 ('cycle:10', {5:'1/5',9:'1/5',10:'1/5',18:'1/5',20:'1/5'},
  ['0','2/5','2/5','0','0','2/5','2/5','0','2/5','0']),
 ('bp:11113', {7:'3/20',11:'3/20',13:'3/20',14:'3/20',16:'2/5'},
  ['3/10','3/10','3/10','0','3/10','3/10','0','3/10','0','0']),
 ('bp:01112', {6:'2/9',10:'2/9',12:'2/9',17:'1/3'},
  ['0','2/9','0','1/3','2/9','2/9','0','2/9','0','0']),
 ('bp:11122', {7:'1/5',8:'1/10',9:'1/10',10:'1/10',12:'1/10',16:'1/10',17:'1/10',18:'1/10',20:'1/10'},
  ['1/5','1/5','1/10','1/10','1/5','1/10','1/10','1/10','1/10','0']),
 ('bp:11111', {1:'1/20',2:'1/20',3:'3/40',4:'1/20',5:'3/40',6:'3/40',8:'1/20',9:'3/40',10:'3/40',12:'3/40',16:'1/20',17:'3/40',18:'3/40',20:'3/40',24:'3/40'},
  ['3/40']*10),
 ('rooted:0', {1:'5/26',6:'3/26',8:'3/26',10:'3/26',12:'3/26',16:'3/26',18:'3/26',20:'3/26'},
  ['0']*4+['3/26']*6),
 ('corrected:11223', {7:'1/6',11:'1/6',12:'1/4',16:'1/12',17:'1/6',18:'1/6'},
  ['1/3','1/6','1/6','1/6','1/6','1/6','1/6','1/4','0','0']),
]

CERTS = make_certificates()
assert len(CERTS) == 98
# Interior reference distribution: five independent Bernoulli(1/5) errors,
# with pair upper bounds exactly 1/25.
ref_pi = [F(1,5)**sum(int(t) for t in x)*F(4,5)**(5-sum(int(t) for t in x)) for x in STATES]
ref_b = [F(1,25)]*10
alpha=F(1,100)

def evaluate(c,q,b):
    return c.a+sum(v*t for v,t in zip(c.beta,q))+sum(v*t for v,t in zip(c.gamma,b))

def verify(label, mass, bounds):
    pi=[F(mass.get(mask, '0')) for mask in range(32)]
    b=[F(u) for u in bounds]
    assert sum(pi)==1 and all(v>=0 for v in pi)
    q0=[sum(pi[mask]*int(STATES[mask,i]) for mask in range(32)) for i in range(5)]
    assert all(b[e]>=sum(pi[mask]*int(STATES[mask,i])*int(STATES[mask,j]) for mask in range(32)) for e,(i,j) in enumerate(EDGES))
    assert all(b[e]<=min(q0[i],q0[j]) for e,(i,j) in enumerate(EDGES))
    pi=[(1-alpha)*a+alpha*t for a,t in zip(pi,ref_pi)]
    b=[(1-alpha)*a+alpha*t for a,t in zip(b,ref_b)]
    q=[sum(pi[mask]*int(STATES[mask,i]) for mask in range(32)) for i in range(5)]
    assert sum(pi)==1 and all(F(0)<t<F(1,2) for t in q)
    assert sum(q)-3*min(q)<=1
    for e,(i,j) in enumerate(EDGES):
        joint=sum(pi[mask]*int(STATES[mask,i])*int(STATES[mask,j]) for mask in range(32))
        assert joint <= b[e] <= min(q[i],q[j])
    k=next(i for i,c in enumerate(CERTS) if c.label==label)
    scores=[evaluate(c,q,b) for c in CERTS]
    gap=min(v-scores[k] for j,v in enumerate(scores) if j!=k)
    assert gap>0
    return gap

if __name__=='__main__':
    for label,mass,b in WITNESSES:
        print(label, 'strict margin =', verify(label,mass,b))
    print('Nine exact rational witnesses + relabeling cover all 98 terms.')
    print('This proves irredundancy within U98, not minimax completeness.')
\end{lstlisting}
\section{1,882 条规则和 60 个置换轨道的整数枚举}
\noindent 保存为 \texttt{D2\_K5\_orbit\_count\_exact.py} 并执行。该脚本不使用浮点 LP。

\begin{lstlisting}
"""Integer-only check: 1,882 signed threshold rules and 60 S5-orbits."""
from itertools import permutations, product
CANON=[(1,0,0,0,0),(3,1,1,1,1),(2,1,1,1,0),
       (3,2,2,1,1),(1,1,1,0,0),(2,2,1,1,1),(1,1,1,1,1)]
REP=[s for s in range(32) if not(s&1)]
PERMS=list(permutations(range(5)))

def bits_from_w(w):
    f=[]
    for s in range(32):
        dot=sum(w[i]*(1 if s&(1<<i) else -1) for i in range(5))
        assert dot!=0
        f.append(int(dot>0))
    assert all(f[s]+f[s^31]==1 for s in range(32))
    return sum(f[s]<<j for j,s in enumerate(REP))

def permute_code(code,perm):
    f=[0]*32
    for j,s in enumerate(REP):
        f[s]=(code>>j)&1
        f[s^31]=1-f[s]
    out=[0]*32
    for s in range(32):
        image=sum(((s>>i)&1)<<perm[i] for i in range(5))
        out[image]=f[s]
    return sum(out[s]<<j for j,s in enumerate(REP))

rules=set()
for w0 in CANON:
    for w in set(permutations(w0)):
        for sg in product((-1,1), repeat=5):
            rules.add(bits_from_w(tuple(w[i]*sg[i] for i in range(5))))
assert len(rules)==1882
seen=set(); orbits=0
for code in rules:
    if code in seen: continue
    orb={permute_code(code,p) for p in PERMS}
    assert orb.issubset(rules)
    seen|=orb
    orbits+=1
assert seen==rules and orbits==60
print('signed self-dual threshold rules:',len(rules),'S5 orbits:',orbits)
\end{lstlisting}
\end{document}
